\subsection{COEFFICIENTS OF A REPRESENTATION}

Let U be an associative algebra with unit element over K, U* the dual of the vector space U and \(\rho\) a representation of U on a vector space E. For \(e \in E\) and \(e' \in E^*\) , let \(\theta(e, e') \in U^*\) be the corresponding coefficient of \(\rho\) (Algebra, Chapter VIII, §13, no. 3). Recall that \(\theta(e, e')(x) = \langle \rho(x)e, e' \rangle\) and that the mapping \(e \mapsto \theta(e, e')\) is for fixed \(e'\) a homomorphism of the U-module E into the U-module U* of the coregular representation of U (loc. cit., Proposition 1); therefore the vector subspace of U* generated by the coefficients of \(\rho\) (a subspace which we shall denote by C( \(\rho\) ) in this paragraph) is a sub-U-module of U*. If \((e_i')_{i \in I}\) is a family of elements generating E* over K, the mapping \(e \mapsto (\theta(e, e_i'))\) is an injective U-homomorphism of E into C( \(\rho\) ) \(^{I}\) , for \(\theta(e, e_i') = 0\) for all i implies \(\langle e, e_i' \rangle = \langle \rho(1)e, e_i' \rangle = \theta(e, e_i')(1) = 0\) for all i and hence e = 0.

In particular, if U is the enveloping algebra of a Lie algebra g and \(\rho\) is a representation of g (identified with a representation of U) on an n-dimensional vector space E, the g-module E is isomorphic to a sub-g-module of \((\mathrm{C}(\rho))^{n}\) .

